\chapter{System Overview}\label{ch:system}
\section{Stakeholders and use context}
The prototype supports a human inspecting an authorized X-ray image. The
operator needs visible task identity, localized findings, usable review
controls and honest handling of missing inputs and failed execution. Students
and supervisors also need reproducible artifacts and enough provenance to
explain implementation decisions. Physical scanner integration requires a
separate observed interface handoff; public documentation and image replay do
not establish compatibility with the lab machine.

The research dataset and the lab collection serve different purposes. Public
IEDXray supports the single-scan baseline and executed generic training/evaluation.
Dismantled scans support the distributed collection workstream. Their counts
do not supply detector labels, independent cases or accepted negatives.

\section{Functional requirements}
\begin{longtable}{p{0.1\textwidth}p{0.49\textwidth}p{0.31\textwidth}}
\caption{Requirement-to-evidence map for the SDP I prototype.}\label{tab:requirements}\\
\toprule ID & Requirement & Evidence and disposition\\\midrule\endfirsthead
\toprule ID & Requirement & Evidence and disposition\\\midrule\endhead
FR1 & Accept authorized images and identify the source explicitly & Upload, folder receipt and finite replay documented; physical scanner not qualified\\
FR2 & Apply declared preprocessing consistently before inference & Canonical image and backend transforms documented; no measured enhancement benefit\\
FR3 & Detect and localize the selected task with a verified checkpoint & Generic published-test measurements in Chapter 5; device/specific diagnostic scope remains bounded\\
FR4 & Display original-coordinate boxes, scores and task identity & Shared contract and GUI implementation evidence\\
FR5 & Preserve execution failure and missing-input states & Contract/software behavior; empty results imply no benign verdict\\
FR6 & Support operator review and export traceable records & Notes, session review and JSON export; persistence beyond session remains limited\\
FR7 & Evaluate with matching ground truth and frozen settings & Generic test job 820 verified; duplicate/group limitations retained\\
FR8 & Associate distributed scans using justified case provenance & Contract direction only; case labels and fusion unfinished\\
\bottomrule
\end{longtable}

The proposal's acquisition objective remains broader than the present inventory.
The recorded catalogue plans 500 images across 125 four-bag groups. The inventory
contains 478 distinct byte payloads; the numerical difference is not an accepted
shortfall until the lab confirms round mapping and usable-count criteria.

\section{System specification}
\subsection{Inputs and outputs}
A successful scan result contains version, scan ID, execution status, original
image dimensions, model/task identity and provenance, threshold, and a detection
list. Each detection carries an ID, original-pixel corner coordinates, a category
namespace and label, a confidence value, and explicit host association where
available. An unmatched region retains no invented host link.

The run manifest binds checkpoint hash, inference-code identity, task maps,
input identities, backend settings and actual execution device. Training settings
and study records remain separate from the inference manifest. This separation
allows a candidate to be evaluated without extending every consumer with
training-specific fields.
\projectrecord{docs/contracts.md; docs/baseline-run-contract.md}

\subsection{Failure and empty-result semantics}
Execution success and detection count are separate. A successful run can contain
zero detections; it means the model reported no regions at the chosen settings.
A failed or missing scan is not an empty successful scan. Application messages
and downstream evaluation must retain this distinction. It is particularly
important for distributed cases, where a missing package cannot be treated as
evidence that a component is absent.

\subsection{Nonfunctional requirements}
Reproducibility requires matching revisions, hashes, dependency profiles and
documented settings. Local processing reduces unnecessary transfer of sensitive
imagery; raw scans, credentials and private manifests stay outside Git. Usability
requires readable overlays, keyboard review, responsive layouts and explicit
source/result identity. The team must verify these properties on the actual
presentation machine before claiming demo readiness.

Latency has at least two meanings: model execution time and total time from
accepted input to a visible completed result. Hardware, precision, batch size,
warm-up, synchronization and preprocessing/postprocessing boundaries must be
reported before comparing speeds. No stable percentile latency is inferred
from a few diagnostic samples.

\section{Acceptance boundaries}
The existing vertical slice is accepted only to its demonstrated extent.
Full P1 requires an agreed host/threat association or cascade and decision policy.
P2 additionally requires component-level evidence, justified cases and case-level
ground truth. Neither is supplied automatically by choosing a checkpoint in the
GUI. Table~\ref{tab:requirements} therefore records partial and unfinished
requirements rather than presenting all boxes as complete.

\pending{Formal supervisor acceptance of the current requirements and
mapping, current presentation-machine qualification and live scanner interface.}
